% Comprehensive, accessible report for SIMC 2026 Question 3 (parts a-h).
% Compile with XeLaTeX (uses the system Times New Roman font), TWO passes for the
% table of contents and cross-references.  Run from the repository root:
%   xelatex -output-directory=doc doc/Q3_full_report.tex   (x2)
\documentclass[11pt]{article}

\usepackage[a4paper,margin=2.3cm]{geometry}
\usepackage{amsmath,amssymb,amsthm}
\usepackage{fontspec}
\setmainfont{Times New Roman}
\usepackage{graphicx}
\usepackage{booktabs}
\usepackage{xcolor}
\usepackage[hidelinks]{hyperref}
\usepackage[font=small,labelfont=bf]{caption}
\usepackage{enumitem}
\setlist{nosep,leftmargin=1.4em}

\graphicspath{{doc/FINAL/figs/}{FINAL/figs/}{figs/}}
\setlength{\parskip}{3pt}
\renewcommand{\arraystretch}{1.15}

\DeclareMathOperator{\tr}{tr}
\DeclareMathOperator{\Var}{Var}
\DeclareMathOperator{\Cov}{Cov}
\newcommand{\R}{\mathbb{R}}
\newcommand{\Xt}{\widetilde{X}}
\newcommand{\xt}{\widetilde{x}}
\newcommand{\Nt}{\widetilde{N}}
\newcommand{\norm}[1]{\lVert #1 \rVert}
\newcommand{\T}{^{\mathsf{T}}}

\theoremstyle{definition}\newtheorem{definition}{Definition}[section]
\theoremstyle{plain}\newtheorem{proposition}{Proposition}[section]
\theoremstyle{remark}\newtheorem{remark}{Remark}[section]

\title{\bfseries Question 3 --- A Complete Report\\[4pt]
\large Cleaning, Principal Component Analysis, and Clustering\\ of a High-Dimensional Spectroscopic Dataset}
\author{SIMC 2026 Endeavour Challenge \quad (Parts a--h)}
\date{May 2026}

\begin{document}
\maketitle

\begin{abstract}\noindent
This report works through the whole of Question 3 in detail and from the ground up. We are
given a table of $1000$ infrared spectra, each recorded at $3401$ frequencies, and asked to
clean it, analyse it, and ultimately recover how many distinct food samples it contains. Along
the way we meet --- and carefully define --- the row norm, mean-centering, variance, the
covariance and Gram matrices, eigenvalues and eigenvectors, Principal Component Analysis (PCA),
the singular value decomposition, cosine similarity, and clustering. The headline results are:
$40$ blank vials and $36$ anomalous high-norm vials are removed (leaving $924$ clean spectra);
the data is strongly low-dimensional, with a dominance ratio $\lambda_1/\lambda_2\approx3.53$
and a clear ``elbow'' at $k^\star=8$ signal directions; the leading principal direction is an
interpretable contrast between molecular bands; and clustering --- carried out two independent
ways that we prove must agree --- shows the genuine vials fall into well-separated groups of
about $37$ repeats each, giving a final count of \textbf{26 unique food samples}.
\end{abstract}

\tableofcontents
\bigskip

% ======================================================================
\section{Introduction and roadmap}\label{sec:intro}

Imagine a blind taste test. A thousand small vials are filled with fermented, caffeinated
drinks --- kombucha, cocoa drinks, coffee liqueurs, Chinese tea --- and an instrument shines
infrared light of many different frequencies through each one, recording how strongly the
liquid absorbs at each frequency. The result is a long list of numbers per vial, a
\emph{spectrum}, which acts as a chemical fingerprint. The labels saying which vial held which
drink have been lost, and our job is to use linear algebra to reconstruct the structure of the
experiment: which measurements are corrupt, which frequencies carry information, what the
dominant patterns of variation are, and --- finally --- how many genuinely different drinks
were tested.

The mathematics needed is introduced as we go, so no prior linear algebra is assumed. The
report is organised as follows. Section~\ref{sec:background} defines every tool we will use.
Section~\ref{sec:data} describes the dataset and a subtlety about how it is stored. Then each
part of Question 3 has its own section: cleaning (\S\ref{sec:3a}), variance and the covariance
matrix (\S\ref{sec:3b}), the eigenspectrum and PCA (\S\ref{sec:3c}), the leading eigenvector
(\S\ref{sec:3d}), clustering in two different spaces (\S\ref{sec:3e}--\ref{sec:3f}), the
comparison of those two approaches (\S\ref{sec:3g}), and the final count of unique samples
(\S\ref{sec:3h}). Section~\ref{sec:conclusion} ties the threads together, and the appendices
give a glossary and a table of all key numbers.

% ======================================================================

\end{document}
